\documentclass[10pt,a4paper]{amsart}
\usepackage[margin=19mm]{geometry}
\usepackage{amsmath,amssymb,mathtools}
\usepackage[scr=rsfs,cal=euler]{mathalfa}
\usepackage{xfrac}
\usepackage[unicode,hidelinks,bookmarks=false]{hyperref}
\newcommand{\ie}{i.e.\ }
\newcommand{\cf}{cf.\ }
\newcommand{\resp}{respectively\ }
\newcommand{\Subsection}{\subsection}
\providecommand{\nobreakdash}{\nobreak\hbox{-}}
\setlength{\emergencystretch}{2em}
\multlinegap0pt
\newtheorem{prop}{Proposition}
\newtheorem{lema}{Lemma}
\newtheorem{thm}{Theorem}
\title{Eigenvalue homogenization for the $p-$Laplacian with rapidly oscillating $L^q$ weights}
\author{Ariel Salort}
\date{Source: arXiv:2608.16231v1 (CC BY 4.0)}
\begin{document}
\maketitle

\noindent\emph{Source and licence.} Eigenvalue homogenization for the $p-$Laplacian with rapidly oscillating $L^q$ weights. Version arXiv:2608.16231v1 (CC BY 4.0), distributed by arXiv under the Creative Commons Attribution 4.0 International licence, http://creativecommons.org/licenses/by/4.0/, as recorded in the arXiv licence metadata for this exact version and shown on the abstract page https://arxiv.org/abs/2608.16231v1. Full licence notice: \texttt{licenses/arxiv-2608.16231-CC-BY-4.0.txt}.\par\smallskip

\noindent\textit{Editorial scope.} Counted unit: the article's thm teo.d (source0.tex lines 687--701). The article's complete original proof of it is delivered with it (source0.tex lines 704--785, one \texttt{proof} environment of 82 lines). No mathematics is written, repaired or re-derived: the statement and the proof are the article's own lines, in the article's own notation, and no retained line is substituted or re-flowed.  Retained uncounted as the source-local context the proof needs: lines 75--80; lines 89--94; lines 146--149; lines 418--432; lines 651--657.  Nothing was omitted from any retained span, so the package carries no omission record.  No retained line contains a citation, so the wrapper carries no bibliography and no bibliography entry.  No retained line contains a figure environment, an image-inclusion command or drawing code, so no image is delivered and the package declares no asset dependency.  Editorial formatting only, in addition: an AMS \texttt{amsart} preamble that restates the article's own declarations and macros used by the retained lines, namely the environments \texttt{prop}, \texttt{lema}, \texttt{thm} on separate counters, together with the standard packages those lines need.  The retained environments print in this document as Proposition 1, Lema 1, Theorem 1 in order of appearance, whereas the article prints the same items with the numbering of its own full document.  The complete native source of the article as distributed in the arXiv e-print for this exact version is bundled unchanged as \texttt{sources/207814/source0.tex}.
\begin{align}\tag{$\mathcal{D}_\varepsilon$}   \label{eq.p.eps} 
\begin{cases}
-\Delta_p u_\varepsilon + V_\varepsilon |u|_\varepsilon^{p-2} u_\varepsilon= \lambda_\varepsilon \rho_\varepsilon |u_\varepsilon|^{p-2}u_\varepsilon &\text{ in } \Omega\\
u_\varepsilon=0 & \text{ on } \partial\Omega.
\end{cases}
\end{align}

\begin{align}\tag{$\mathcal{D}_0$}   \label{eq.p.0} 
\begin{cases}
-\Delta_p u_0 +  V_0 |u_0|^{p-2}u_0 = \lambda_0 \rho_0 |u_0|^{p-2}u_0 &\text{ in } \Omega\\
u_0=0 & \text{ on } \partial\Omega.
\end{cases}
\end{align}

\begin{align} 
\frac{n}{p} < q\leq \infty &\quad \text{ when }1<p\leq n  \label{rel1}\\
q=1 &\quad \text{ when } p>n \label{rel2}
\end{align}

\begin{prop} \label{prop.dirichlet}
Let $p$ and $q$ satisfying condition \eqref{rel1}--\eqref{rel2}. Then
$$
\left|\int_\Omega(\rho_\varepsilon-\bar\rho) |u|^p\,dx \right|  \leq   C \varepsilon^\alpha \|\rho_\varepsilon - \bar \rho\|_{L^q(\Omega)}   \|\nabla u\|_{L^p(\Omega)}^p
$$
holds for all $u\in W^{1,p}_0(\Omega)$ and $\varepsilon\ll 1$, where 
\begin{align} \label{alfa}
\alpha=
\begin{cases}
p-\frac{n}{q} & \text{ when } q> \frac{n}{p} \text{ and } 1<p\leq n\\
1-\frac{n}{p} & \text{ when } q= 1 \text{ and } p>n
\end{cases}
\end{align}
and $C$ is a constant depending only of $p$  and $\Omega$.
\end{prop}

\begin{lema} \label{cota.rho}
Let $1\leq q\leq \infty$. Given a bounded open set $\Omega\subset \mathbb R^n$,  we have that
$$
\|\rho_\varepsilon - \bar \rho\|_{L^q(\Omega)} \leq 
\|\rho\|_{L^q(Q)}\left( \mathop{diam}(\Omega)^\frac{n}{q} + |\Omega|^\frac1q \right).
$$
\end{lema}

\begin{thm} \label{teo.d}
  Let $\lambda_{k,\varepsilon}$ and $\lambda_{k,0}$ the $k-$th eigenvalue of problems \eqref{eq.p.eps} and \eqref{eq.p.0}, respectively. Then there are (computable) constants $C$ and $\textbf{C}$ depending only of $p$, $q$, $\rho$, $V$ and $\Omega$ such that
\begin{align*}
|\lambda_{k,0}-\lambda_{k,\varepsilon}| &\leq  
C  \varepsilon^\alpha \left(\lambda_{k,0} \|V_\varepsilon - \bar V\|_{L^q(\Omega)}  +  \lambda_{k,0}^2 \|\rho_\varepsilon - \bar \rho\|_{L^q(\Omega)}  \right)
\end{align*}
holds for $0<\varepsilon<(2\textbf{C} \lambda_{k,0})^{-\frac{1}{\alpha}}$, where
$$
\alpha=
\begin{cases}
p-\frac{n}{q} & \text{ when } q> \frac{n}{p} \text{ and } 1<p\leq n\\
1-\frac{n}{p} & \text{ when } q=1 \text{ and } p>n.
\end{cases}
$$
\end{thm}

\begin{proof}
Let $\delta>0$ and let $G_\delta^k\subset W^{1,p}_0(\Omega)$ be a compact, symmetric set of genus $k$ such that 
$$
\lambda_{k,0} = \inf_{G\in \Gamma_k} \sup_{u\in G} \frac{\int_\Omega |\nabla u|^p + \bar V |u|^p \,dx}{\int_\Omega \bar \rho |u|^p\,dx} 
= \sup_{u\in G_\delta^k} \frac{\int_\Omega |\nabla u|^p + \bar V |u|^p\,dx}{\int_\Omega \bar \rho |u|^p\,dx} + O(\delta).
$$
The set $G_\delta^k$ is admissible in the variational characterization of $\lambda_{k,\varepsilon}$, then
\begin{align*}
\lambda_{k,\varepsilon} &\leq \sup_{u\in G_\delta^k} \frac{\int_\Omega |\nabla u|^p  +   V_\varepsilon |u|^p\,dx}{\int_\Omega  \rho_\varepsilon |u|^p\,dx}\\ 
&=  \sup_{u\in G_\delta^k} \left( \frac{\int_\Omega |\nabla u|^p + \bar V |u|^p \,dx}{\int_\Omega  \bar \rho |u|^p\,dx}  +  \frac{\int_\Omega (V_\varepsilon - \bar V) |u|^p \,dx}{\int_\Omega  \bar \rho |u|^p\,dx} \right)
\frac{\int_\Omega  \bar \rho |u|^p\,dx}{\int_\Omega  \rho_\varepsilon |u|^p\,dx} . 
\end{align*}
Observe that, for every $u\in G_\delta^k$ we have that
\begin{equation} \label{eqd.1}
\frac{\int_\Omega |\nabla u|^p + \bar V |u|^p \,dx}{\int_\Omega  \bar \rho |u|^p  \,dx} \leq  \sup_{u\in G_\delta^k} \frac{\int_\Omega |\nabla u|^p + \bar V |u|^p\,dx}{\int_\Omega  \bar \rho |u|^p\,dx} = \lambda_{k,0}+O(\delta). 
\end{equation}
Proposition \ref{prop.dirichlet} and \eqref{eqd.1} yield
\begin{align} \label{eqd.2}
\begin{split}
\frac{\int_\Omega (V_\varepsilon - \bar V) |u|^p \,dx}{\int_\Omega  \bar \rho |u|^p\,dx} &\leq  C_{V_\varepsilon}
 \varepsilon^\alpha    \frac{\int_\Omega  |\nabla u|^p \,dx}{\int_\Omega  \bar \rho |u|^p\,dx}\leq
 C_{V_\varepsilon} \varepsilon^\alpha   (\lambda_{k,0}+O(\delta)),
\end{split}
\end{align}
where we have denoted $C_{V_\varepsilon}:=C\|V_\varepsilon - \bar V\|_{L^q(\Omega)}$.


Since $u\in G_\delta^k\subset W^{1,p}_0(\Omega)$, by using Proposition \ref{prop.dirichlet}, \eqref{eqd.1} and Lemma \ref{cota.rho} we obtain that
\begin{align*}
 \frac{\int_\Omega  \bar \rho |u|^p\,dx}{\int_\Omega  \rho_\varepsilon |u|^p\,dx} &\leq 1 + C_{\rho_\varepsilon} \varepsilon^\alpha  \frac{\int_\Omega |\nabla u|^p\,dx}{\int_\Omega \rho_\varepsilon |u|^p\,dx}\\
&= 1 + C_{\rho_\varepsilon} \varepsilon^\alpha   \frac{\int_\Omega |\nabla u|^p\,dx}{\int_\Omega \bar \rho |u|^p\,dx} \frac{\int_\Omega \bar \rho |u|^p\,dx}{\int_\Omega \rho_\varepsilon |u|^p\,dx}\\
&\leq 1 + C_\rho \varepsilon^\alpha (\lambda_{k,0}+O(\delta)) \frac{\int_\Omega \bar \rho |u|^p\,dx}{\int_\Omega \rho_\varepsilon |u|^p\,dx}
\end{align*}
where $C_\rho$ is the bound of $C_{\rho_\varepsilon}$ (independent of $\varepsilon$ and $k$) given in  Lemma \ref{cota.rho}. Then, for $\varepsilon$ small enough, i.e., $\varepsilon<(C_\rho ( \lambda_{k,0}+O(\delta))^{-\frac{1}{\alpha}}$, this gives that
$$
\frac{\int_\Omega  \bar \rho |u|^p\,dx}{\int_\Omega  \rho_\varepsilon |u|^p\,dx}\leq \frac{1}{1- C_\rho \varepsilon^\alpha  (\lambda_{k,0}+O(\delta))}.
$$
Hence, the last two inequalities yield
\begin{equation} \label{eqd.3}
\frac{\int_\Omega  \bar \rho |u|^p\,dx}{\int_\Omega  \rho_\varepsilon |u|^p\,dx}\leq 1+ \varepsilon^\alpha \frac{ C_{\rho_\varepsilon} (\lambda_{k,0}+O(\delta))}{1- C_\rho\varepsilon^\alpha  (\lambda_{k,0}+O(\delta)) }.
\end{equation}
Gathering \eqref{eqd.1}, \eqref{eqd.2} and \eqref{eqd.3} and letting $\delta\to 0$, we get
\begin{align*}
\lambda_{k,\varepsilon}&\leq \lambda_{k,0}(1 +C_{V_\varepsilon} \varepsilon^\alpha  ) \left( 1+ \varepsilon^\alpha\frac{ C_{\rho_\varepsilon}   \lambda_{k,0}}{1- C_\rho \varepsilon^\alpha   \lambda_{k,0} } \right)
\end{align*}
that is, 
\begin{equation} \label{eq.r.1}
\lambda_{k,\varepsilon}-\lambda_{k,0} \leq C_{V_\varepsilon} \lambda_{k,0} \varepsilon^\alpha +  \frac{2 C_{\rho_\varepsilon} (1+C_{V_\varepsilon}) \lambda_{k,0}^2}{1- C_\rho  \varepsilon^\alpha \lambda_{k,0}}\varepsilon^\alpha
\end{equation}
and it holds for $\varepsilon< \min\{1,( C_\rho\lambda_{k,0})^{-\frac{1}{\alpha}}\}$.


We will obtain now an inequality similar to \eqref{eq.r.1} which bounds the difference $\lambda_{k,0}-\lambda_{k,\varepsilon}$. For this end, we observe first that   $\lambda_{k,\varepsilon}$ can be bounded independently of $\varepsilon$. For instance, if we take  $\varepsilon=  (2  C_\rho \lambda_{k,0})^{-\frac{1}{\alpha}}$ we can obtain from \eqref{eq.r.1} the following:
\begin{align*}
\lambda_{k,\varepsilon} &\leq \lambda_{k,0} +  C_V \lambda_{k,0}\varepsilon^\alpha + \frac{2C_\rho (1+C_V)\lambda_{k,0}^2}{1-C_\rho \varepsilon^\alpha \lambda_{k,0}}\varepsilon^\alpha\leq \lambda_{k,0} +  \frac{  C_V}{2C_\rho}  +    2 (1+C_V) \lambda_{k,0}
\end{align*}
(where the constant $C_V$ is defined as  $C_\rho$ and it is independent of $\varepsilon$), then
\begin{equation} \label{desig1}
\lambda_{k,\varepsilon} \leq   \left(3+ \frac{ C_V }{2 C_\rho}\frac{1}{\lambda_{1,0}} +2 C_V \right) \lambda_{k,0} = \tilde C \lambda_{k,0},
\end{equation}
where $\tilde C=\tilde C(p,V,\rho,\Omega)>1$ is independent of $\varepsilon$ and $k$.

Using \eqref{desig1} and a reasoning similar to the one leading to  \eqref{eq.r.1}, we can obtain that
\begin{align} \label{eq.r.2}
\begin{split}
\lambda_{k,0}-\lambda_{k,\varepsilon} \leq\tilde C C_{V_\varepsilon} \lambda_{k,0} \varepsilon^\alpha +   \frac{2 \tilde C^2  C_{\rho_\varepsilon} (1+C_{V_\varepsilon}) \lambda_{k,0}^2}{1- C_\rho \tilde C  \varepsilon^\alpha \lambda_{k,0}} \varepsilon^\alpha
\end{split}
\end{align}
for $\varepsilon <(\bar C_\rho \tilde C \lambda_{k,0})^{-\frac{1}{\alpha}}$. From \eqref{eq.r.1} and \eqref{eq.r.2} we obtain that
\begin{align*}
|\lambda_{k,0}-\lambda_{k,\varepsilon}| &\leq  
\tilde C C_{V_\varepsilon} \lambda_{k,0} \varepsilon^\alpha +  \frac{2 \tilde C^2 C_{\rho_\varepsilon} (1+C_{V_\varepsilon}) \lambda_{k,0}^2}{1-C_\rho \tilde C \varepsilon^\alpha \lambda_{k,0}}\varepsilon^\alpha 
\end{align*}
for $\varepsilon < \min\{1,(\bar C_\rho \tilde C \lambda_{k,0})^{-\frac{1}{\alpha}}\}$.

Finally, by using again Lemma \ref{cota.rho}, for some constants $C$ and $\textbf{C}$  depending only of $p$,  $V$, $\rho$ and $\Omega)$ we can rewrite the previous relation as
\begin{align*}
|\lambda_{k,0}-\lambda_{k,\varepsilon}| &\leq  
C  \varepsilon^\alpha \|V_\varepsilon - \bar V\|_{L^q(\Omega)} \lambda_{k,0} + C  \|\rho_\varepsilon - \bar \rho\|_{L^q(\Omega)}  \varepsilon^\alpha \frac{ \lambda_{k,0}^2}{1- \textbf{C} \varepsilon^\alpha \lambda_{k,0}}
\end{align*}
and it holds for instance when $\varepsilon < (2\textbf{C} \lambda_{k,0})^{-\frac{1}{\alpha}}$. This completes the proof.
\end{proof}

\end{document}
